% HD-R18. GPT-6 Astra Pro /root; inherited effort not exposed.
% 21 November 2026. The original primitive rank-four coefficient branch.
\section{HD-R18: the actual primitive rank-four core}
\label{sec:rank-four-primitive-core}

Retain the original natural finite-curve source for a still-missing
Weil-sixfold determinant component, over
\(K=\mathbb Q(i)\) or \(K=\mathbb Q(\sqrt{-d})\),
\(d=2,7,11,19,43,67,163\). Its actual faithful coefficient
group is \(G\subset GL_m(K)\), its parabolic coefficient
cohomology has dimension six, and its period differential has
rank nine. In the remaining branch considered here, every
nonidentity element of \(G\) has moving rank at least four.
Let \(N=N_4\) be the actual normal subgroup generated by its
moving-rank-four elements. The branch in which \(N\) is
absolutely irreducible and complex primitive satisfies
\[
                         21\le m\le63.                 \tag{R18.1}
\]
It has an actual characteristic quasi-simple subgroup
\(S\triangleleft N\), absolutely irreducible on the original
\(K^m\). With \(Q=G\cap K^*I\), the original conjugation maps give
\[
 \begin{gathered}
 G/Q\lhook\joinrel\longrightarrow\operatorname{Aut}(S/Z(S)),
 \qquad\operatorname{Inn}(S/Z(S))\subset\operatorname{image}(G/Q),\\
                  |Z(S)|\mid\gcd(m,|\mu(K)|).           \tag{R18.2}
 \end{gathered}
\]
Combining this with the proved field bounds in HD-R12 gives
\(21\le m\le63\) for \(d=43\), \(33\le m\le63\) for
\(d=67\), and excludes this primitive branch for \(d=163\).
The finite actual models and their normalizing matrices in the
remaining ranges still require analysis; no exclusion of those
ranges is asserted in this section.

\paragraph{1. The degree bound on the original matrices.}
The actual sphere Euler equation and the marked-base parameter
bound are
\[
                  6=-2m+\sum_{j=1}^b c_j,
            \qquad b\ge12,\qquad c_j\ge4.               \tag{R18.3}
\]
The lower bound in (R18.1) follows. HD-R17 applies to \(N\),
since \(m\ge21>16\); it excludes every nontrivial tensor
structure, including permutations of tensor factors.
Any block system or tensor structure preserved by \(G\)
would also be preserved by \(N\). Thus \(G\) is primitive
and tensor indecomposable. The irreducible group \(N\) is
nontrivial, so at least one actual rank-four generator exists.
Consequently \(\nu_G(V)\le4\), in the notation of HD-G31.

Use the exact characteristic-zero statement and the complete
exception list of Guralnick--Saxl, Theorems 7.1 and 7.A,
already read and recorded in HD-G31. Outside the scalar
deleted-permutation alternative and the listed degrees at
most ten, its strict inequality is
\(\nu_G(V)>\max\{2,\sqrt m/2\}\).
For \(m\ge64\) this forces the scalar deleted-permutation
alternative. It supplies an actual complex linear isomorphism
\(C_0:V_{\mathbb C}\to V_0\), where
\[
 V_0=\{z\in\mathbb C^s:\sum z_i=0\},\quad s=m+1,
 \qquad C_0\rho(g)C_0^{-1}=a_gP_{\pi(g)},
 \qquad\pi(G)=A_s\ \hbox{or}\ S_s.                     \tag{R18.4}
\]
The scalar and permutation factors are unique: a permutation
scalar on the augmentation space preserves each line
\(\mathbb C(e_i-e_j)\), hence fixes every unordered pair
and every label for \(s\ge3\). Its scalar is then one.
Thus \(\pi\) is an actual homomorphism.

Choose original lifts \(g,h\in G\) of
\(u=(12)(45)\) and \(v=(23)(45)\). The scalar factors cancel,
and the original commutator satisfies
\[
 C_0\rho([g,h])C_0^{-1}
                   =[P_u,P_v]=P_{(132)}.                \tag{R18.5}
\]
The permutation has moving rank two on \(V_0\): its fixed
space on \(\mathbb C^s\) has dimension \(s-2\), and
removing the constant line leaves fixed dimension \(s-3\)
in dimension \(s-1\). Rank is unchanged on extending
the original \(K\)-matrix to \(\mathbb C\).
This gives an actual forbidden rank-two element of \(G\).
It proves \(m\le63\), including exclusion of the endpoint64
by the strict published inequality. This contradiction uses
the literal commutator, with no replacement of the source.

\paragraph{2. Normal restrictions and the elementary-abelian case.}
Put \(Z=Z(N)=N\cap K^*I\). Every nonscalar normal subgroup
\(H\triangleleft N\) is absolutely irreducible. Indeed its
complex isotypic summands are permuted transitively by \(N\);
primitivity makes their number one. For its unique irreducible
type \(W\), the actual evaluation map is
\[
 W\otimes\operatorname{Hom}_H(W,V_{\mathbb C})
        \xrightarrow{\sim}V_{\mathbb C},\qquad
                       w\otimes f\longmapsto f(w).       \tag{R18.6}
\]
The image algebra is \(\operatorname{End}(W)\otimes I\).
An original element of \(N\) normalizing this algebra
acts on its matrix units by an inner implementation;
after removing it, the element commutes with all matrix
units and acts on the multiplicity space. Hence (R18.6)
is an actual fixed tensor structure for \(N\). HD-R17
forces multiplicity one, since \(\dim W=1\) would make
\(H\) scalar. Thus \(Z(H)\) consists of scalar matrices;
for the subgroups \(H\) containing \(Z\) used below,
this gives \(Z(H)=Z\).

Choose a minimal nontrivial normal subgroup \(M\) of
\(N/Z\), and let \(H\) be its full inverse image. The
elementary minimal-normal-subgroup argument proved in
HD-G32 gives either \(M=C_p^b\) or \(M=T^t\), with
\(T\) nonabelian simple. In the first case the actual
commutators give a nondegenerate alternating pairing
\[
 H/Z\times H/Z\longrightarrow\mu_p(K),\qquad
                        (\bar u,\bar v)\longmapsto[u,v].
                                                               \tag{R18.7}
\]
The radical is zero because \(Z(H)=Z\); its values are
\(p\)-torsion because \(u^p\in Z\). The pairing has
a nonidentity value since an irreducible abelian group
has dimension one. The eight fields displayed above
have no odd-prime roots of unity, so \(p=2\).
Successive nonzero paired vectors split off symplectic
planes, giving \(H/Z=\mathbb F_2^{2a}\).

For each coset \(v\), choose an original representative
and its matrix \(B_v\in GL_m(K)\), with \(B_0=I\).
Conjugation by a representative paired nontrivially with
\(v\ne0\) sends \(B_v\) to \(-B_v\); thus
\(\operatorname{tr}B_v=0\) and
\[
 \operatorname{tr}(B_u^{-1}B_v)=m\delta_{u,v}.
                                                               \tag{R18.8}
\]
The averaging operator on \(\operatorname{End}(V)\)
projects onto its one-dimensional scalar commutant.
Taking its trace, and using (R18.8) on the scalar cosets,
gives
\[
 1=|H|^{-1}\sum_{h\in H}|\operatorname{tr}\rho(h)|^2
       =m^2/|H/Z|,\qquad m=2^a.                         \tag{R18.9}
\]
Therefore the original \(B_v\) form a \(K\)-basis of
\(\operatorname{End}_K(K^m)\). For any original \(x\in N\),
conjugation has the literal form
\[
 xB_vx^{-1}=\lambda_vB_{Av},\qquad
 A\in\operatorname{Aut}_{\mathbb F_2}(H/Z),\quad
 \lambda_v\in\mu(K).                                   \tag{R18.10}
\]
On \(F=\ker(A-I)\), the multiplication identities
\(B_uB_v=c(u,v)B_{u+v}\), with original central scalar
\(c(u,v)\), show that \(\lambda\) is a character of \(F\).
Taking the trace of (R18.10) on the full matrix algebra gives
\[
 |\operatorname{tr}x|^2
       =\sum_{v\in F}\lambda_v\in\{0,2^{\dim F}\}.
                                                               \tag{R18.11}
\]
The equality with the squared absolute trace uses finite
order: every eigenvalue has modulus one, so the trace of
the inverse is the complex conjugate trace.

The only power of two in \([21,63]\) is32. Take an actual
rank-four generator \(x\). It has28 eigenvalues equal to
one and four other roots of unity. Hence
\(\operatorname{Re}\operatorname{tr}x\ge24\), while
\(|\operatorname{tr}x|<32\). The latter inequality is
strict: equality in the triangle inequality would force
all32 unit eigenvalues to have the same argument; the28
ones would then force the other four to equal one.
Consequently
\[
                  576\le|\operatorname{tr}x|^2<1024.      \tag{R18.12}
\]
This interval contains neither zero nor a power of two,
since \(512<576\) and the next power is1024.
It contradicts (R18.11). The elementary-abelian case is
therefore excluded directly in the original matrix algebra.

\paragraph{3. The perfect factors and their literal tensor map.}
We now have \(H/Z=T^t\). Put \(E=H'\). Perfection of
\(H/Z\) gives \(H=EZ\), and then
\(E=[EZ,EZ]=[E,E]\). Set \(Z_E=E\cap Z\), let
\(H_i\) be the inverse image of \(T_i\) in \(E\),
and put \(E_i=H_i'\). Again \(H_i=E_iZ_E\), and each
\(E_i\) is perfect. For \(i\ne j\), commutators between
\(H_i\) and \(H_j\) are central. Commuting with a fixed
element of \(H_j\) is therefore a homomorphism
\(H_i\to Z_E\), which kills \(E_i\). Thus the \(E_i\)
commute; their product has central quotient in the perfect
\(E\), and
\[
        E=E_1\cdots E_t,\qquad
        E_i/Z(E_i)\simeq T_i.                           \tag{R18.13}
\]
A central element of \(E_i\) commutes with all factors
and therefore with \(H=EZ\); it is scalar by absolute
irreducibility. This verifies the centre assertion in
(R18.13), so every \(E_i\) is quasi-simple.

The restriction to each \(E_i\) has one isotypic type:
its isotypic projectors commute with all factors and hence
with the irreducible \(E\). Choose matrix units
\(e^{(i)}_{ab}\) in their commuting simple image algebras.
Products of diagonal units sum to \(I\), and products of
off-diagonal units identify the resulting joint corners.
The corner \(\prod_i e^{(i)}_{11}V_{\mathbb C}\) is
nonzero and one-dimensional. A larger corner would carry
a nonscalar operator extending through those same matrix
units to a nonscalar commutant of \(E\). Choose its
nonzero vector \(u\). The actual tensor isomorphism is
\[
 \bigotimes_i e_{a_i}\longmapsto
            \left(\prod_i e^{(i)}_{a_i1}\right)u.         \tag{R18.14}
\]
Its inverse records the coefficient of \(u\) in
\(\prod_i e^{(i)}_{1a_i}v\) for every multi-index.
The matrix-unit identities verify both composites.
Each factor dimension exceeds one, since a nontrivial
perfect factor cannot have a faithful nontrivial scalar
image. Conjugation by \(N\) permutes the actual \(E_i\)
and their algebras, so it preserves (R18.14) with its
factor permutations. HD-R17 forces \(t=1\).
Thus \(S=E\) is quasi-simple and absolutely irreducible
on the original representation.

\paragraph{4. Characteristic core, field maps and the same corner.}
The minimal normal subgroup \(M\) of \(N/Z\) is unique.
A distinct one has full inverse image \(H_2\) satisfying
\([H_2,H]\subset Z\). The central-commutator homomorphism
kills the perfect subgroup \(S=H'\), giving
\([H_2,S]=1\). Its absolute irreducibility then makes
\(H_2\) scalar, contradicting nontriviality of that
minimal subgroup. Hence \(H\), and then \(S=H'\), are
characteristic in \(N\). The original \(G\) normalizes
\(N\) because conjugation preserves moving rank; it
therefore normalizes \(S\).

Conjugation \(G\to\operatorname{Aut}(S)\) has kernel
exactly \(Q\). An automorphism of the perfect \(S\)
trivial on \(S/Z(S)\) has the form \(s\mapsto sz_s\),
where \(s\mapsto z_s\) is a homomorphism to the centre.
It vanishes by perfection. This proves the injection in
(R18.2), whose image contains inner automorphisms from
the actual subgroup \(S\). Its centre is contained in
\(\mu(K)I\); the determinant character of the perfect
\(S\) is trivial, so each central scalar satisfies
\(z^m=1\). This proves the divisibility in (R18.2).

For every induced automorphism \(\alpha\), the normalizing
matrices remain the solutions of the original finite
\(K\)-linear equations
\[
          T\rho(s)=\rho(\alpha(s))T\qquad(s\in S).       \tag{R18.15}
\]
An actual generating set suffices. The solution space
contains the original normalizing matrix and has complex
dimension one by absolute irreducibility. Gaussian
elimination over \(K\) has the same rank after extension,
so (R18.15) is a one-dimensional \(K\)-space and every
nonzero solution is invertible. These equations retain
the original outer-intertwiner lines and all scalar phases.

Throughout the argument \(S\), \(N\), and \(G\) are
actual subgroups of the original \(GL_m(K)\); the cover,
primitive vector, rational matrix-coefficient idempotent,
integer-cleared Jacobian image and restricted polarization
are unchanged. A later \(K\)-coordinate map \(C\) sends
every original matrix to \(C\rho(g)C^{-1}\) and the
primitive line \(Kw\) to \(KCw\). With a prescribed
positive invariant comparison form \(h_0\), its exact
relation to the original form is
\(h_0(Cx,Cy)=a h(x,y)\), \(a\in\mathbb Q_{>0}\).
The normalized coefficient is consequently identical:
\[
 \frac{h_0(Cw,C\rho(g)w)}{h_0(Cw,Cw)}
                 =\frac{h(w,\rho(g)w)}{h(w,w)}.          \tag{R18.16}
\]
The Schur factors satisfy \(c_0=c/a\), hence
\(c_0a=c\) on the same paired factor. When changing a
primitive line, the actual Fourier matrix unit of HD-G31
satisfies \(b^\dagger b=(q_v/q_w)e_w\); its integer-cleared
isogeny has polarization multiplier \(M_0^2q_v/q_w\).
These are the precise coefficient and Jacobian maps available
for the remaining finite model analysis.
